\documentclass[a4paper,12pt]{article}
\usepackage[utf8]{inputenc}
\usepackage[T1]{fontenc}
\usepackage[ngerman]{babel}
\usepackage{geometry}
\usepackage{hyperref}
\usepackage{amsmath}
\usepackage{booktabs}
\usepackage{parskip}
\geometry{margin=2.5cm}

\begin{document}

\section*{Physics-Constrained Deep Learning for High-dimensional Surrogate Modeling
and Uncertainty Quantification without Labeled Data}

\textbf{Autoren:} Yinhao Zhu, Nicholas Zabaras, Phaedon-Stelios Koutsourelakis, Paris Perdikaris \\
\textbf{Journal:} Journal of Computational Physics \\
\textbf{Jahr:} 2019 (arXiv: 1901.06314)

\subsection*{Motivation und Problemstellung}

Surrogat-Modelle für PDE-Systeme werden üblicherweise als überwachte Lernprobleme formuliert,
bei denen Paare aus Eingabe- und Ausgabedaten für das Training benötigt werden.
Das grundlegende Problem dabei ist, dass die Erzeugung dieser Ausgabedaten (d.\,h. das
Lösen der Gleichungen mit klassischen Simulatoren) rechenintensiv und oft der eigentliche
Flaschenhals in der wissenschaftlichen Praxis ist. Die Autoren identifizieren zwei zentrale
Herausforderungen: (1) den hohen Datenbedarf tiefer neuronaler Netze im Vergleich zur
begrenzten Verfügbarkeit teurer Simulationsdaten (``small data problem''), und (2) die
fehlende Berücksichtigung physikalischer Nebenbedingungen in rein datengetriebenen Ansätzen.

\subsection*{Kernbeitrag: Physik-beschränktes Surrogat-Training}

Das Hauptmerkmal der Arbeit ist ein Framework, das die physikalischen Gleichungen direkt
in die Verlustfunktion integriert und damit das Training \emph{ohne} gelabelte
Ausgabedaten ermöglicht. Statt der klassischen mittleren quadratischen Abweichung zwischen
Vorhersage und Referenzsimulation wird die sogenannte \emph{Physics-Constrained Loss} verwendet:
%
\begin{equation}
\mathcal{L}(\theta; \{x^{(i)}\}) = \frac{1}{N} \sum_{i=1}^{N} \bigl[ V(\hat{y}_\theta(x^{(i)}), x^{(i)}) + \lambda\, B(\hat{y}_\theta(x^{(i)})) \bigr],
\end{equation}
%
wobei $V(\cdot)$ der PDE-Verlustterm (z.\,B. das Residuum oder ein Variationsfunktional)
und $B(\cdot)$ der Randbedingungs-Verlustterm ist. Das Netz lernt also, die Gleichung
zu lösen, anstatt eine vorberechnete Lösung nachzuahmen.

Die Arbeit untersucht vier Varianten der Verlustfunktion für Darcy-Strömungen:
\emph{Primal Residual}, \emph{Primal Variational}, \emph{Mixed Variational} und
\emph{Mixed Residual}. Die mixed-residual Formulierung wird bevorzugt, da sie die
Differentiationsordnung reduziert (approximiert durch Sobel-Filter) und die beste
Vorhersagegenauigkeit für das Druckfeld liefert.

\subsection*{Architektur: Convolutional Encoder-Decoder}

Als Surrogat-Architektur verwenden die Autoren ein \emph{Dense Convolutional Encoder-Decoder}
Netz (ähnlich dem U-Net~\cite{ronneberger2015unet}), das aus drei dichten Blöcken mit je 6,
8 und 6 Schichten (Growth Rate 16) aufgebaut ist. Das Netz empfängt ein einzelnes
Eingabekanal-Bild der Permeabilitätsfeld-Realisierung ($1 \times 64 \times 64$) und gibt
drei Ausgabekanäle aus: Druckfeld und zwei Flussfeld-Komponenten ($3 \times 64 \times 64$).
Räumliche Gradienten werden durch Sobel-Filter approxi\-miert, was gegenüber automatischer
Differentiation durch das FC-NN deutlich effizienter ist.

Ein zentrales Ergebnis ist die Überlegenheit von CNNs gegenüber vollständig verbundenen
Netzen (FC-NNs) für diesen Aufgabentyp: CNNs konvergieren deutlich schneller, erfassen
multi\-skalige Strukturen der Strömungsfelder besser und scheitern nicht bei komplexen
Eingabefeldern wie kanalisierten Strukturen, bei denen FC-NNs keine ausreichende Genauigkeit
erreichen.

\subsection*{Probabilistisches Surrogat und Unsicherheitsquantifizierung}

Neben dem deterministischen Surrogat wird ein probabilistisches Modell vorgestellt,
das auf einem konditionierten \emph{Flow-basierten generativen Modell} (Conditional Glow)
basiert. Das Modell wird durch Minimierung der \emph{reversen Kullback-Leibler-Divergenz}
zwischen der Modellverteilung $p_\theta(\mathbf{y}|\mathbf{x})$ und einer
Boltzmann-Gibbs-Referenzverteilung trainiert:
%
\begin{equation}
\text{D}_{\mathrm{KL}}(p_\theta(\mathbf{y}|\mathbf{x}) \| p_\beta(\mathbf{y}|\mathbf{x})) = \beta\, \mathbb{E}[\mathcal{L}(\mathbf{y},\mathbf{x})] - H_\theta(\mathbf{y}|\mathbf{x}) + \text{const},
\end{equation}
%
wobei $\beta$ ein inverser Temperatur-Parameter ist, der die Breite der Vorhersageverteilung
kontrolliert. Der erste Term erzwingt die Erfüllung der PDE, der zweite Term (negative
Entropie) fördert Vielfalt und verhindert Modenkollaps. Das Modell ermöglicht exakte
Log-Likelihood-Bewertung und generiert Unsicherheitsabschätzungen ohne Ausgabedaten.

\subsection*{Experimentelle Ergebnisse}

Die Methode wird an stationären Darcy-Strömungen in zufälligen heterogenen Medien
evaluiert. Als Eingaben dienen fünf Verteilungstypen: Gausssche Zufallsfelder (GRF)
mit unterschiedlicher KLE-Trunkierung (KLE128 bis KLE2048), verzerrte GRFs und
kanalisierte Felder. Wichtige quantitative Resultate:

\begin{itemize}
    \item Das physik-beschränkte Surrogat (PCS) erzielt für das Druckfeld relative
          $L_2$-Fehler von ca. 1,7\,\% (bei 8192 Trainingsdaten, GRF KLE512),
          vergleichbar mit dem datengetriebenen Surrogat (DDS) bei 2,0\,\%.

    \item Bei der Generalisierung auf \emph{Out-of-Distribution}-Eingaben (kanalisierte
          Felder, verzerrte GRFs) zeigt das PCS konsistent bessere Generalisierungs\-fähigkeit
          als das DDS, da die Physik als Regularisierung wirkt.

    \item Das probabilistische Surrogat (Conditional Glow, $\beta = 150$) erreicht
          einen relativen $L_2$-Fehler von 0,38\,\% für das Druckfeld und gut
          kalibrierte Unsicherheitsabschätzungen.
\end{itemize}

\subsection*{Bezug zur Masterarbeit}

Die Arbeit von Zhu et al.\ berührt mehrere Kernthemen der vorliegenden Masterarbeit
und bietet sowohl konzeptionelle Inspiration als auch wichtige Vergleichspunkte:

\textbf{1. Convolutional Encoder-Decoder als Surrogat für Strömungsfelder.}
Die verwendete Netzarchitektur ist konzeptionell direkt verwandt mit der in dieser
Masterarbeit eingesetzten U-Net-Architektur. Beide verwenden Skip-Connections zwischen
Encoder- und Decoder-Pfad, mehrskalige Feature-Extraktion und lernen eine
Abbildung von einem skalaren Eingabefeld auf vektorielle Strömungsfelder.
Ein wesentlicher Unterschied besteht darin, dass das Eingabefeld bei Zhu et al.\ das
Permeabilitätsfeld darstellt, während in der Masterarbeit die Positionen und
Geometrien der sedi\-men\-tierenden Kristalle als Eingabe dienen.

\textbf{2. Physik-informierte Verlustterme.}
Die Integration physikalischer Nebenbedingungen als Verlustterme -- insbesondere
Kontinuitätsgleichungen und Divergenzfreiheit -- ist auch für die Vorhersage von
Stokes-Strömungsfeldern um Kristalle relevant. Die Masterarbeit implementiert eine
ähnliche Strategie mit einem Kontinuitätsverlust $\mathcal{L}_{\text{cont}} = \|\nabla \cdot
\mathbf{v}\|^2$, der die Inkompressibilität des Fluids erzwingt. Das in Zhu et al.\
beschriebene Residual-Learning mit gemischter Formulierung bietet dabei Anknüpfungs\-punkte
für alternative Verlustformulierungen im Stokes-Regime.

\textbf{3. Residual-Learning-Ansatz.}
Obwohl Zhu et al.\ keinen expliziten Residual-Learning-Ansatz im Sinne von
$\mathbf{v}_{\text{total}} = \mathbf{v}_{\text{Stokes}} + \Delta\mathbf{v}_{\text{gelernt}}$
verfolgen, entspricht die mixed-residual Verlustfunktion konzeptionell dem Prinzip,
das Netz auf die Abweichung von einer analytischen Basislösung zu trainieren.
Die Masterarbeit erweitert dieses Prinzip durch das explizite Einbetten der analytischen
Stokes-Lösung als Priori, was eine stärkere physikalische Regularisierung bietet.

\textbf{4. Generalisierung auf unbekannte Konfigurationen.}
Ein Kernbefund von Zhu et al.\ ist, dass physik-beschränkte Modelle bessere
Out-of-Distribution-Generalisierung zeigen als rein datengetriebene Methoden.
Dieser Befund ist direkt relevant für die zentrale Forschungsfrage der Masterarbeit:
Kann ein Netz, das auf $N$ Kristallen trainiert wurde, auf $N{+}m$ Kristalle generalisieren?
Die Ergebnisse von Zhu et al.\ stützen die Hypothese, dass Physik-informierung die
Generalisierungsfähigkeit systematisch verbessert -- insbesondere wenn die Testverteilung
von der Trainingsverteilung abweicht, wie es beim Übergang von $N$- zu $(N{+}m)$-Kristall-
Konfigurationen der Fall ist.

\textbf{5. Vergleich CNN vs.\ FC-NN.}
Der klare Befund, dass CNNs FC-NNs für räumliche Strömungsfelder überlegen sind,
unterstützt die Architekturwahl der Masterarbeit. Insbesondere die effizientere
Erfassung multi\-skaliger Strukturen -- die im Stokes-Regime durch
Wechselwirkungen zwischen mehreren Kristallen auf verschiedenen Längenskalen
entstehen -- spricht für die U-Net-basierte Herangehensweise.

\textbf{Wesentliche Unterschiede und Erweiterungen.}
Im Gegensatz zu Zhu et al., die ein festes räumliches Medium mit stochastischen
Eigenschaften modellieren, muss die Masterarbeit strukturelle Variabilität in der
Anzahl und Position der Kristalle berücksichtigen. Dies entspricht einem Szenario,
in dem sich nicht nur die PDE-Koeffizienten, sondern auch die Geometrie des Problems
verändert -- eine deutlich anspruchsvollere Generalisierungsaufgabe. Darüber hinaus
konzentriert sich die Masterarbeit auf das Stokes-Regime (Re $\ll 1$) mit bekannter
analytischer Einzellösungsstruktur, während Zhu et al.\ allgemeine Darcy-Strömungen
behandeln.

\subsection*{Einschätzung und Relevanz}

Die Arbeit von Zhu et al.\ ist eine der ersten, die physik-beschränktes Training für
hochdimensionale PDE-Surrogat-Modelle systematisch demonstriert. Sie liefert starke
empirische Evidenz dafür, dass die Integration physikalischer Nebenbedingungen sowohl
die Dateneffizienz als auch die Generalisierungsfähigkeit verbessert. Für die
Masterarbeit ist die Arbeit als methodische Referenz für physik-informierte Encoder-Decoder-
Architekturen besonders wertvoll -- insbesondere hinsichtlich der Verlustformulierung und
der Bewertung von Generalisierungsleistung auf Out-of-Distribution-Szenarien.

\begin{thebibliography}{9}
\bibitem{ronneberger2015unet}
O.~Ronneberger, P.~Fischer, T.~Brox,
\emph{U-Net: Convolutional Networks for Biomedical Image Segmentation},
MICCAI 2015.
\end{thebibliography}

\end{document}